
\documentclass[11pt]{article}

\usepackage[margin=1in]{geometry}
\usepackage{amsmath,amssymb,amsfonts}
\usepackage{booktabs}
\usepackage[hidelinks]{hyperref}

\title{Representation and Time--Memory Tradeoffs in Symmetric Cryptography:\\
Finite-Field Basis Transformations for AES and Reproducible Hellman Experiments}
\author{Ali Mkhida\\Algorizk Labs\\[2pt]\small Current draft maintainer\\\small Originating from an M2P SCCI master's research project at Ensimag--UJF\\\small developed with Abdourahmane Sakho and Maad El Yadari}
\date{Working ePrint draft --- 28 August 2026}

\newcommand{\F}{\mathbb{F}}
\newcommand{\GF}{\mathbb{F}}
\newcommand{\Tr}{\operatorname{Tr}}
\newcommand{\wt}{\operatorname{wt}}
\newcommand{\Enc}{\operatorname{Enc}}
\newcommand{\Inv}{\operatorname{Inv}}
\newcommand{\im}{\operatorname{im}}

\begin{document}
\maketitle

\begin{abstract}
This note upgrades a historical master's research project on high-performance
cryptography into a reproducible research artifact.  The original presentation
combined finite-field representation changes for AES-like arithmetic with
practical secret-key recovery experiments based on exhaustive search, Hellman
time--memory tradeoffs, and distinguished points.  The repository now preserves
the original presentation as a historical PDF, while this note gives a separate
ePrint-style reconstruction.

The first part reconstructs the representation mathematics directly from the
preserved slides: the concrete $\GF(16)$ normal basis, its exact bilinear
multiplication formula, the squaring permutation, the displayed 8-by-8
isomorphism matrices, and the two-level composite-field non-LUT inversion
proposal.  The second part gives a modern reproducible
implementation of the cryptanalytic experiments over deliberately reduced DES
state spaces.  The main empirical observation is that, for the tested
single-map Hellman construction with $mt/N=1$, the median exact table coverage
is only about $16.8\%$, far below the independent-sample occupancy reference
$1-e^{-1}\simeq 63.2\%$.  The gap is explained by functional-graph coalescence:
chain states are not independent samples but iterates of one fixed map.
\end{abstract}

\section{Historical artifact and updated goal}

The original project was titled \emph{High Performance Computing for
Cryptography}.  That title correctly described the original research context, but it does
not reveal the mathematical object studied by the modernized repository.  The
updated repository is organized around two related themes:

\begin{enumerate}
  \item finite-field representation changes for efficient AES-like arithmetic;
  \item cryptanalytic time--memory tradeoffs for reduced key spaces.
\end{enumerate}

The historical presentation remains preserved unchanged in the repository at
\[
  \texttt{docs/original-presentation.pdf}.
\]
This document is separate: it is the research-note version, intended to evolve
toward an ePrint-style manuscript.
Where the original slides were visual and compressed, this note makes the
algebraic objects explicit and separates historical claims from newly
reproduced measurements.

The reconstruction below is based directly on slides 4--19 of the preserved
presentation~\cite{historical}.  Slides 20 and 21 contain only the phrase ``PROOF ON BOARD'' for,
respectively, secret-key recovery and the claimed 200~MHz AES throughput.
Consequently, this note does not reconstruct or endorse a proof or a
200~MHz implementation result from those two slides; such claims require
independent source material or a new implementation.

\section{Finite fields in characteristic two}

Let $n\geq 1$.  The binary extension field $\GF(2^n)$ may be represented as
\[
  \GF(2^n) \simeq \GF(2)[X]/(p(X)),
\]
where $p$ is an irreducible polynomial of degree $n$.  If $\alpha$ denotes the
class of $X$, the \emph{polynomial basis} is
\[
  (1,\alpha,\alpha^2,\ldots,\alpha^{n-1}).
\]
Every field element is written uniquely as
\[
  a = a_0 + a_1\alpha + \cdots + a_{n-1}\alpha^{n-1},
  \qquad a_i\in\GF(2).
\]

A second important representation is a \emph{normal basis}.  An element
$\beta\in\GF(2^n)$ is normal over $\GF(2)$ if
\[
  \mathcal N(\beta)=
  (\beta,\beta^2,\beta^{2^2},\ldots,\beta^{2^{n-1}})
\]
is a basis of $\GF(2^n)$ over $\GF(2)$.  In this representation, the Frobenius
map
\[
  \varphi : x \longmapsto x^2
\]
acts as a cyclic shift on coordinates:
\[
  \sum_{i=0}^{n-1} b_i \beta^{2^i}
  \longmapsto
  \sum_{i=0}^{n-1} b_i \beta^{2^{i+1}}.
\]
This is the central reason normal bases are attractive in hardware:
squaring can be implemented by rewiring rather than by a general multiplier.


\section{The historical \texorpdfstring{$\GF(16)$}{GF(16)} model}

The historical presentation does not work only at the abstract
$\GF(2^n)$ level.  It fixes the concrete quartic field
\[
  K_4
  =
  \GF(2)[X]/(X^4+X^3+X^2+X+1)
\]
and lets $\alpha$ denote the residue class of $X$.  The slides first use the
ordered polynomial (``canonical'') basis
\[
  \mathcal B_{\mathrm{can}}
  =
  (\alpha^3,\alpha^2,\alpha,1).
\]
Because
\[
  \alpha^4+\alpha^3+\alpha^2+\alpha+1=0,
\]
we have $\alpha^5=1$ and $\alpha\neq 1$.

The presentation then uses the ordered normal basis
\[
  \mathcal B_{\mathrm{nor}}
  =
  (\alpha^8,\alpha^4,\alpha^2,\alpha).
\]
The displayed identities are
\[
  1_{K_4}=\alpha^8+\alpha^4+\alpha^2+\alpha,
\]
together with
\[
  \alpha^3=\alpha^8,\qquad
  \alpha^5=\alpha^{10}
    =\alpha^8+\alpha^4+\alpha^2+\alpha,
\]
\[
  \alpha^6=\alpha^{16}=\alpha,\qquad
  \alpha^9=\alpha^4,\qquad
  \alpha^{12}=\alpha^2.
\]
These are consistent with $\alpha^5=1$.  In particular, Frobenius squaring
cycles the basis elements
\[
  \alpha^8\mapsto\alpha,\qquad
  \alpha^4\mapsto\alpha^8,\qquad
  \alpha^2\mapsto\alpha^4,\qquad
  \alpha\mapsto\alpha^2.
\]

For coordinate vectors
\[
  x=[a_3,a_2,a_1,a_0]_{\mathcal B_{\mathrm{nor}}},
  \qquad
  y=[b_3,b_2,b_1,b_0]_{\mathcal B_{\mathrm{nor}}},
\]
this means
\[
  x=a_3\alpha^8+a_2\alpha^4+a_1\alpha^2+a_0\alpha,
\]
and similarly for $y$.

\subsection{Exact multiplication formula from the slides}

Slide~7 expands the product in the normal basis.  Defining the common term
\[
  q
  =
  a_0b_2\oplus a_2b_0\oplus a_1b_3\oplus a_3b_1,
\]
the four output coordinates can be written compactly as
\begin{align*}
  c_3 &=
  a_0b_1\oplus a_1b_0\oplus a_2b_2\oplus q,\\
  c_2 &=
  a_0b_3\oplus a_3b_0\oplus a_1b_1\oplus q,\\
  c_1 &=
  a_2b_3\oplus a_3b_2\oplus a_0b_0\oplus q,\\
  c_0 &=
  a_1b_2\oplus a_2b_1\oplus a_3b_3\oplus q.
\end{align*}
Thus
\[
  xy
  =
  c_3\alpha^8+c_2\alpha^4+c_1\alpha^2+c_0\alpha.
\]
The repeated term $q$ is precisely the common subexpression highlighted in the
historical slide.  Sharing it is a direct logic optimization: the same XOR
subnetwork feeds all four output coordinates.

The repository now includes a small verifier that exhaustively checks this
formula on all $16^2$ input pairs against multiplication modulo
$X^4+X^3+X^2+X+1$.

\subsection{Squaring is a coordinate permutation}

Setting $y=x$ in the displayed multiplication formula, or simply applying the
Frobenius map, gives the historical slide-8 identity
\[
  x^2
  =
  a_2\alpha^8+a_1\alpha^4+a_0\alpha^2+a_3\alpha.
\]
Hence
\[
  [a_3,a_2,a_1,a_0]
  \longmapsto
  [a_2,a_1,a_0,a_3].
\]
No field multiplier is required for this operation in the chosen normal basis:
squaring is a cyclic coordinate permutation.  This is the precise content of
the slide's statement that squaring ``comes for free in hardware.''  The same
normal-basis principle underlies inversion algorithms such as
Itoh--Tsujii~\cite{itoh1988}.

\section{The tower-field isomorphism used for the S-box}

AES represents a byte as an element of
\[
  K_8
  =
  \GF(2)[T]/(T^8+T^4+T^3+T+1),
\]
the standard AES field~\cite{fips197}.  The historical presentation writes the
S-box schematically as
\[
  S(t)=A_s t^{-1}+B_s,
\]
where the subscript here only disambiguates the affine-map constants from the
quadratic-extension coefficients used below.

The central proposal in slide~5 is to choose a field isomorphism
\[
  \delta:
  K_8
  \longrightarrow
  K_4[Y]/(Y^2+A_qY+B_q)
\]
and perform the multiplicative inversion in the quadratic extension of
$K_4$:
\[
  t
  \xrightarrow{\delta}
  a_1Y+a_0
  \xrightarrow{\mathrm{inverse}}
  b_1Y+b_0
  \xrightarrow{\delta^{-1}}
  t^{-1}.
\]
This is the exact representation-level decomposition shown in the
presentation.  It is closely related to the composite-field AES S-box
architectures studied by Satoh et al.~\cite{satoh2001} and
Canright~\cite{canright2005}.

\subsection{Quadratic-extension inversion}

Let
\[
  z=a_1Y+a_0,
  \qquad
  Y^2+A_qY+B_q=0.
\]
The conjugate root of $Y$ is $Y+A_q$, so
\[
  \bar z
  =
  a_1Y+(a_0+A_q a_1).
\]
Its norm down to $K_4$ is
\[
  N(z)
  =
  z\bar z
  =
  a_0^2+A_q a_0a_1+B_q a_1^2.
\]
For $z\neq 0$,
\[
  z^{-1}
  =
  \frac{a_1Y+(a_0+A_q a_1)}
       {a_0^2+A_q a_0a_1+B_q a_1^2}.
\]
The normal basis from the previous section makes both squarings in the
denominator coordinate permutations.  This explains why the historical
project first optimized $K_4$ arithmetic and only then returned to the
8-bit S-box.

\section{Constructing and representing the isomorphism}

The slides describe the isomorphism by matching generators.  Let
$K_8^\times$ and
\[
  L^\times
  =
  \left(K_4[Y]/(Y^2+A_qY+B_q)\right)^\times
\]
be the multiplicative groups.  If $\alpha_8\in K_8^\times$ and
$\beta\in L^\times$ are roots of the same irreducible degree-8 polynomial
$P$, define
\[
  \delta(\alpha_8)=\beta,
  \qquad
  \delta(\alpha_8^i)=\beta^i.
\]
Because
\[
  |K_8^\times|=2^8-1=255=3\cdot 5\cdot 17,
\]
a candidate $g$ is primitive exactly when
\[
  g^{255/3}\neq 1,\qquad
  g^{255/5}\neq 1,\qquad
  g^{255/17}\neq 1.
\]
This is the standard precise generator test corresponding to the abbreviated
criterion written on slide~11.

The historical slide states the candidate
\[
  \beta=[1,0,0,1]\,Y+[0,0,0,0]
\]
under its $K_4$ coordinate convention and records
\[
  P(\beta)=0,
  \qquad
  P(T)=T^8+T^4+T^3+T+1.
\]
It then defines a derived map $\delta_2$ by its image of the polynomial
generator and extends it multiplicatively:
\[
  \delta_2(T^i)=\delta_2(T)^i.
\]
The target coefficients $A_q,B_q$ are not written on the surviving slides,
so the root calculation itself cannot be independently reconstructed from the
PDF alone.  What \emph{is} recoverable exactly is the resulting binary
linear map.

\subsection{Exact \texorpdfstring{$8\times 8$}{8x8} matrices}

Slide~12 gives the matrix of $\delta_2$ and its inverse under the bit ordering
used in the presentation:
\[
D=
\left[
\begin{array}{cccccccc}
1&0&1&0&0&0&1&0\\
0&1&1&1&0&0&0&0\\
1&1&0&1&1&1&0&0\\
1&0&1&0&1&1&1&0\\
1&0&1&0&0&1&0&1\\
1&0&1&0&0&0&0&1\\
0&1&0&1&1&1&0&1\\
0&1&0&0&0&0&0&1
\end{array}
\right],
\]
and
\[
D^{-1}=
\left[
\begin{array}{cccccccc}
1&1&0&1&0&1&1&0\\
1&1&1&1&0&1&0&1\\
0&0&1&0&0&1&1&0\\
1&0&0&1&0&0&1&1\\
1&0&0&1&1&1&0&0\\
0&0&0&0&1&1&0&0\\
0&1&1&1&0&0&0&0\\
1&1&1&1&0&1&0&0
\end{array}
\right].
\]
The modernization verifier checks over $\GF(2)$ that
\[
  DD^{-1}=D^{-1}D=I_8.
\]

\section{Hardware meaning of the isomorphism matrices}

For an 8-bit coordinate vector $v$, the basis conversion is simply
\[
  \delta_2(v)=Dv,
  \qquad
  \delta_2^{-1}(v)=D^{-1}v
  \qquad\text{over }\GF(2).
\]
Each output bit is therefore an XOR of selected input bits, and all eight
output coordinates can be computed in parallel.  This is exactly the point
made in slides~14--15.

A natural structural cost is the matrix Hamming weight
\[
  \operatorname{wt}(M)=
  |\{(i,j):M_{ij}=1\}|.
\]
For the matrices printed in the historical presentation,
\[
  \operatorname{wt}(D)=30,
  \qquad
  \operatorname{wt}(D^{-1})=32.
\]
The slide explicitly suggests searching for an isomorphism with low Hamming
weight.  In modern terms, one should optimize not only total ones, but also
shared XOR subexpressions, XOR depth, fanout, and placement/routing delay.
Canright's later systematic treatment of basis choices and isomorphism
matrices is directly relevant to this optimization viewpoint~\cite{canright2005}.

\section{A second composite-field layer and the non-LUT proposal}

Slides~17--19 propose a second decomposition to shorten the longest
combinational path.  The historical notation uses maps $\delta_a$ and
$\delta_b$ and states the goal of factoring the first isomorphism through a
smaller field.  To avoid the composition-order ambiguity in the slide
notation, we write the two operational stages as
\[
  \Delta_8:
  K_8\longrightarrow K_4[Y]/(Y^2+A_qY+B_q)
\]
and
\[
  \Delta_4:
  K_4\longrightarrow
  \GF(4)[Z]/(Z^2+CZ+D_q).
\]
Extend $\Delta_4$ coefficientwise to the quadratic extension.  The combined
representation change is then unambiguously
\[
  \Psi=\widehat{\Delta}_4\circ\Delta_8.
\]

The pipeline corresponding to slide~19 is
\[
  t
  \xrightarrow{\Delta_8}
  a_1Y+a_0
  \xrightarrow{\widehat{\Delta}_4}
  \widetilde a
  \xrightarrow{\mathrm{inverse}}
  \widetilde c
  \xrightarrow{\widehat{\Delta}_4^{-1}}
  (a_1Y+a_0)^{-1}
  \xrightarrow{\Delta_8^{-1}}
  t^{-1}.
\]
The key arithmetic observation is that for every nonzero
$\gamma\in\GF(4)$,
\[
  \gamma^3=1
  \qquad\Longrightarrow\qquad
  \gamma^{-1}=\gamma^2.
\]
Thus inversion in $\GF(4)$ is just squaring on nonzero elements.  In a normal
basis, that squaring is again a coordinate permutation.  This is the precise
mathematical interpretation of the historical statement that the smallest
inversion stage ``comes for free in hardware,'' and it explains the proposed
non-LUT S-box datapath.

The historical slides do not provide the actual $C,D_q$, the concrete
$\Delta_4$ matrix, a synthesized critical path, or an implementation report.
Those are therefore future-work items rather than reconstructed results.


\section{From representation arithmetic to cryptanalytic experiments}

The finite-field part of the project concerns the cost of one encryption
datapath.  The cryptanalytic part concerns the cost of searching over many
candidate keys.  These are linked but distinct:

\begin{itemize}
  \item finite-field basis choices change the cost of encryption rounds;
  \item time--memory tradeoffs change the cost profile of key recovery.
\end{itemize}

The original master's research project connected these views by considering secret-key
recovery experiments after the AES arithmetic discussion.  The modern
repository makes the second part reproducible over deliberately reduced DES
state spaces.  DES is obsolete and is used here only because it provides a
small, standard block cipher setting suitable for reproducible experiments.

\section{Reduced DES state model}

The historical Java prototype varied four raw DES bytes.  However, a DES key
is encoded in eight bytes with one parity bit per byte.  Thus four varying raw
bytes do not automatically give $32$ independent effective DES key bits.

The modern model defines a reduced state space explicitly:
\[
  X_b = \{0,1,\ldots,2^b-1\},
  \qquad 1\le b\le 28.
\]
A state $x\in X_b$ is mapped to a canonical DES key encoding $K(x)$ with odd
parity.  The first four DES bytes are fixed, and the final four bytes carry
up to seven effective bits each.  This gives at most $28$ effective variable
bits.

For a fixed plaintext $P$, the encryption oracle is
\[
  C_x = \Enc_{K(x)}(P).
\]
A reduction function
\[
  R:\{0,1\}^\ast\to X_b
\]
maps ciphertexts back into reduced states.  The core functional graph is
\[
  F(x)=R(\Enc_{K(x)}(P)).
\]
Hellman chains and distinguished-point chains are both walks in this directed
functional graph.

\section{Exhaustive search}

Given a target ciphertext $C^\star$, exhaustive search tests every state:
\[
  \Enc_{K(x)}(P)\stackrel{?}{=}C^\star.
\]
It requires no offline table and no stored endpoints.  Its memory cost is
constant, but its online cost is $O(2^b)$ encryptions in the worst case.

In the verified experiment with $b=16$, exhaustive search recovers all
targets because it covers the entire reduced state space by definition.

\section{Hellman chains}

A Hellman table~\cite{hellman1980} is generated by choosing start states
\[
  s_1,\ldots,s_m\in X_b
\]
and iterating $F$ for a fixed chain length $t$:
\[
  s_i,\ F(s_i),\ F^2(s_i),\ldots,F^t(s_i).
\]
Only the pair
\[
  (s_i,F^t(s_i))
\]
is stored.  The second component is the endpoint.

For a target ciphertext $C^\star$, the online phase computes
\[
  R(C^\star),\quad F(R(C^\star)),\quad F^2(R(C^\star)),\ldots
\]
and checks whether one of these projected endpoints appears in the table.
When an endpoint matches, the corresponding stored chain is regenerated from
its start point and the target ciphertext is verified directly.

The modern implementation stores endpoints in a hash index rather than
repeatedly scanning every endpoint.  It also retains endpoint collisions and
regenerates all candidate chains, because different starts may lead to the
same endpoint.

\subsection{Why the independent occupancy heuristic is optimistic}

If $mt$ states were sampled independently and uniformly from a universe of
size $N=2^b$, the expected covered fraction would be
\[
  1-\left(1-\frac{1}{N}\right)^{mt}
  \approx 1-e^{-mt/N}.
  \label{eq:occupancy}
\]
But a Hellman table does not draw $mt$ independent samples.  It samples start
points, then applies the same function $F$ repeatedly.  If two chains meet,
they coalesce forever:
\[
  F^a(s_i)=F^b(s_j)
  \quad\Longrightarrow\quad
  F^{a+r}(s_i)=F^{b+r}(s_j)
  \quad\text{for all } r\ge 0.
\]
Thus the functional-graph structure of $F$ can reduce actual coverage far
below the independent-sampling reference.

This is not a measurement anomaly; it is a structural effect.  It is also why
exact coverage analysis is included in the repository rather than relying on
the reference value in Eq.~\eqref{eq:occupancy}.

\section{Distinguished points}

A distinguished-point method, closely related to the collision-search framework of van Oorschot and Wiener~\cite{vanoorschot1999}, stores only chains that reach an easily tested
subset $D\subset X_b$.  In the modern implementation,
\[
  D_d=\{x\in X_b : x\bmod 2^d=0\}.
\]
Equivalently, the lowest $d$ state bits are zero.  Under a uniform-state
heuristic, the expected hitting time is approximately $2^d$ transitions.

A chain starts at $s_i$ and is followed until either:

\begin{enumerate}
  \item it reaches a distinguished point; or
  \item it exceeds a maximum chain length.
\end{enumerate}

Only successful chains are stored.  Chains that do not hit a distinguished
point before the cutoff are counted as truncated.  During lookup, the target
is projected forward until it reaches a distinguished point or the bound is
exceeded.  Candidate chains sharing that endpoint are regenerated and verified
against the target ciphertext.

\section{Implementation corrections}

The modern implementation changes the historical prototype in several ways:

\begin{enumerate}
  \item The reduced key space counts effective DES key bits, not raw bytes.
  \item Endpoints are indexed rather than searched by nested linear scans.
  \item Candidate chains are always verified against the target ciphertext.
  \item Distinguished-point chains are bounded and truncated chains are counted.
  \item The implementation records exact state coverage outside the timed path.
  \item All experiments are generated from explicit seeds.
  \item A separate historical-math verifier checks the slide-7 $\GF(16)$
        multiplication formula, the slide-8 squaring permutation, and
        $DD^{-1}=I_8$ for the slide-12 matrices.
\end{enumerate}

These changes make the code a reference implementation rather than a
one-off demonstration.  They also separate algorithmic claims from JVM timing
noise.

\section{Verified experiment}

The committed benchmark uses:

\begin{itemize}
  \item effective state dimension $b=16$, so $N=65\,536$;
  \item $30$ deterministic seeded trials;
  \item Hellman parameters $m=1024$, $t=64$;
  \item distinguished-point parameters $1024$ starts, $d=6$, and maximum
        chain length $256$;
  \item a separate warm-up run before collecting timings;
  \item exact unique-state table coverage computed outside the timed path.
\end{itemize}

\begin{center}
\begin{tabular}{lccc}
\toprule
Method & Exact recovery & Median exact coverage & Median online time \\
\midrule
Exhaustive & $30/30$ & $100.0\%$ & $25.863$ ms \\
Hellman & $4/30$ & $16.8\%$ & $3.399$ ms \\
Distinguished points & $3/30$ & $15.4\%$ & $1.790$ ms \\
\bottomrule
\end{tabular}
\end{center}

For Hellman, the parameters satisfy
\[
  \frac{mt}{N}=\frac{1024\cdot 64}{65\,536}=1.
\]
The independent-sample occupancy reference would therefore be
\[
  1-e^{-1}\simeq 63.2\%.
\]
The exact measured median coverage is instead $16.8\%$.  This confirms that
the fixed-map functional graph is the right mathematical model for these
tables.

For distinguished points, the expected waiting-time reference is
\[
  2^d=64
\]
transitions.  Across the $30$ tables, the median exact stored-table coverage
is $15.4\%$, and many starts merge onto the same distinguished endpoint.

\section{Rust reference implementation and equivalence roadmap}

The Java implementation remains the current reference semantics.  On
28 August 2026 we added an initial Rust reference implementation rather than a
performance rewrite.  The Rust crate reproduces the reduced DES state encoding,
PKCS\#5-padded DES oracle, legacy reduction, exhaustive search, Hellman chains,
distinguished-point chains, and exact coverage analysis.  Shared reference
vectors are consumed by both Java and Rust tests.  Before any cross-language
performance claim, the Rust implementation must continue to reproduce:

\begin{enumerate}
  \item the same reduced DES state encoding;
  \item the same key parity convention;
  \item the same reduction function;
  \item the same chain endpoints for fixed seeds;
  \item the same exact coverage results;
  \item the same recovery decisions on committed test vectors.
\end{enumerate}

Only after cross-language equivalence is established beyond the committed
reference vectors should performance work begin.  The implementation roadmap is:

\[
\begin{aligned}
\text{Java reference}
&\longrightarrow \text{test vectors}
 \longrightarrow \text{Rust reference (initial implementation)}\\
&\longrightarrow \text{equivalence tests}
 \longrightarrow \text{parallel/SIMD optimization}.
\end{aligned}
\]

The committed Rust crate deliberately has no performance claim yet.  Future
work includes property-based testing, Criterion-style benchmarking, larger
cross-language equivalence suites, and parallel chain generation.  A future
hardware path would
treat encryption and reduction as pipeline kernels, but that is outside the
current repository.

\section{Future research directions}

Several directions follow naturally from this update.

\paragraph{Multiple reductions.}
Classical Hellman tables often use table-dependent reductions.  The current
implementation uses one fixed reduction.  Introducing indexed reductions
$R_j$ would reduce chain merging and make the comparison with classical
tradeoff theory more direct.

\paragraph{Rainbow and multi-table variants.}
Rainbow tables change the reduction per column.  A reproducible comparison
between fixed-reduction Hellman tables, multiple Hellman tables, and rainbow
tables would clarify how much of the measured coverage loss is due to
functional-graph coalescence.

\paragraph{Parameter optimization.}
The current benchmark is deliberately small.  A larger study should sweep
$m$, $t$, distinguished-point difficulty $d$, and chain bounds, reporting
coverage, online work, false candidates, and memory.

\paragraph{Finite-field arithmetic implementation.}
The AES finite-field part should be turned into executable code: basis-change
matrices, composite-field inversion, $\delta$ multiplication, and non-LUT
S-box circuits.  The code should report XOR count, depth, and test equivalence
against the AES S-box.

\paragraph{Rust and hardware.}
Rust should first match Java semantics and only then pursue speed.  Hardware
work should be based on stable test vectors and should separate AES arithmetic
optimization from time--memory tradeoff table generation.

\section{Conclusion}

The updated repository separates historical provenance from modern
reproducibility.  The original presentation remains intact, while this note
reconstructs the finite-field mathematics and connects it to a tested
time--memory tradeoff implementation.

The most important experimental lesson is that Hellman table size alone does
not determine coverage.  For the tested fixed reduction, $mt/N=1$ but exact
coverage is only about $16.8\%$.  The correct model is the functional graph of
$F(x)=R(\Enc_{K(x)}(P))$, where chain coalescence is structural.  This makes
the project a useful base for a verified Rust implementation and for a broader
study of reductions, table families, and cryptographic datapath
representations.

\begin{thebibliography}{9}

\bibitem{hellman1980}
M. E. Hellman.
\newblock A cryptanalytic time-memory trade-off.
\newblock \emph{IEEE Transactions on Information Theory}, 26(4):401--406,
1980.

\bibitem{vanoorschot1999}
P. C. van Oorschot and M. J. Wiener.
\newblock Parallel collision search with cryptanalytic applications.
\newblock \emph{Journal of Cryptology}, 12:1--28, 1999.

\bibitem{fips197}
National Institute of Standards and Technology.
\newblock \emph{Advanced Encryption Standard (AES)}, FIPS 197, updated 2023.

\bibitem{rijndael}
J. Daemen and V. Rijmen.
\newblock \emph{The Design of Rijndael: AES -- The Advanced Encryption Standard}.
\newblock Springer, 2002.

\bibitem{canright2005}
D. Canright.
\newblock A very compact S-box for AES.
\newblock In \emph{Cryptographic Hardware and Embedded Systems -- CHES 2005},
pages 441--455. Springer, 2005.

\bibitem{satoh2001}
A. Satoh, S. Morioka, K. Takano, and S. Munetoh.
\newblock A compact Rijndael hardware architecture with S-box optimization.
\newblock In \emph{Advances in Cryptology -- ASIACRYPT 2001}, pages 239--254.
Springer, 2001.


\bibitem{itoh1988}
T. Itoh and S. Tsujii.
\newblock A fast algorithm for computing multiplicative inverses in
$\mathrm{GF}(2^m)$ using normal bases.
\newblock \emph{Information and Computation}, 78(3):171--177, 1988.
\newblock DOI: 10.1016/0890-5401(88)90024-7.

\bibitem{historical}
A. Sakho, A. Mkhida, and M. El Yadari.
\newblock \emph{High Performance Computing for Cryptography}.
\newblock M2P SCCI, Ensimag--UJF, historical master's research presentation, undated.

\end{thebibliography}

\end{document}
